A tabela \ref{tab:limites_juntasPosVelocidade_rad} a seguir complementa a tabela \eqref{tab:limites_juntasPosVelocidade} com os dados convertidos, no caso das juntas revolutas para radianos (rad) e radianos por segundo (rad/s), e para a junta prismática permanecendo em metro (m) e metro por segundo (m/s).

\begin{table}[ht]
\centering
{
\setlength{\tabcolsep}{12pt}
\renewcommand{\arraystretch}{1}
\caption{Limites físicos e operacionais das juntas do manipulador robótico}
\label{tab:limites_juntasPosVelocidade_rad}
\begin{tabular}{clrrrr}
\toprule
\textbf{Junta} & \textbf{Tipo} & $q_{j,\min}$ & $q_{j,\max}$ & $\dot{q}_{j,\min}$ & $\dot{q}_{j,\max}$ \\
\midrule
1 & Revoluta   & $-6{,}2832$ & $6{,}2832$ &
$-0{,}5\pi$ & $0{,}5\pi$ \\[4pt]
2 & Revoluta   & $-1{,}5708$ & $1{,}5708$ &
$-\pi$ & $\pi$ \\[4pt]
3 & Revoluta   & $-2{,}2689$ & $2{,}2689$
& $-\pi$ & $\pi$ \\[4pt]
4 & Revoluta   & $-2{,}6180$ & $2{,}6180$ &
$-2\pi$ & $2\pi$ \\[4pt]
5 & Prismática & $-0{,}15$ m & $0{,}00$ m &
$-0{,}314$ m/s & $0{,}314$ m/s \\[4pt]
\bottomrule
\end{tabular}
}
\end{table}
